\scanpage{600}

\mspage{9}

Supposons à présent que tous les $k_{\unsure{\xi|\zeta}}$ soient égaux au même \deleted{corps commutatif $k$}\inserted{anneau $\boxed{k\ne0}$}. \deleted{Donc les} \deleted{des lors $\Phi$ est un ultrafiltre, $k_\Phi$ est une extension \unsure{de}} \inserted{$k\mathrel{\unsure{-|\to}}k_J\mathrel{-}k_\Phi$} les $k_J$, $k_\Phi$ sont des \deleted{$k$-algèbres,}\inserted{``anneaux sur $k$''}, d'ailleurs $k\mathrel{\unsure{-|\to}}k_{\unsure{\Phi|J}}$ \inserted{(si $J\ne\varnothing$) et $k\mathrel{-}k_\Phi$} est injectif. On se pose la question quand \inserted{$k\mathrel{-}k_\Phi$} c'est \uline{bijectif}, i.e surjectif. Comme $k_\Phi\simeq k_E/\unsure{\mathfrak m|m}_\Phi$, cela signifie donc que $k_E=k+\mathfrak m_\Phi$, i.e que \deleted{de tte} $f\in k_E$ s'écrit

\[
 f=c+g\qquad\left|\begin{aligned}
 &c\in k\quad\text{i.e. }\ c(\xi)=\underset{\substack{\begin{tikzpicture}[baseline=-0.4ex]\node[inner sep=0pt,rotate=90]{$\in$};\end{tikzpicture}\\ k}}{c_0}\quad\forall\xi\in E\\
 &g\in\mathfrak m_\Phi\quad\text{i.e. }\ \exists A\in\Phi\ \text{avec}\\
 &\hspace{9em} g\mid A=0\ .
 \end{aligned}\right.
\]

Mais pour $f\inserted{\in k_E},\ c\in k$ \inserted{et $A\in\Phi$} fixés, l'existence de $g\in \deleted{\unsure{k_\Phi}} k_E$ \deleted{et $A\in\Phi$ comme dessus,} \deleted{tels qu'on ait $f=c+g$, équivaut à l'existence} \inserted{tels que $g\mid A=0$ et $f=c+g$ équivaut} \deleted{de $A\in\Phi$} à \deleted{\unsure{i.e.} $A\subset$} $f\mid A=c_A$, \unsure{(}fonction \unsure{constte|constante} \unsure{sur|de} $A$, de valeur $c$) \inserted{i.e. $A\subset f^{-1}(\{c\})$. Pour $f,c$ fixés, l'existence} \deleted{\unsure{Donc}} de $A\in\Phi$ \deleted{qui \unsure{ait}} \unsure{satisfaisant} ces relations équivaut à $f^{-1}(\{c\})\in\Phi$.

% slanted
\margin{Et pour $f$ fixé, l'exist\textsuperscript{ce} de $c$ équivaut à celle d'une fibre $E_c$ de $f:E\mathrel{\unsure{-|\to}}k$, qui soit dans $\Phi$. D'où}

\uline{Proposition 3} Soit \deleted{$H=\mathrm{Card}\,k$, et} $\Phi\in\unsure{\mathscr F|\mathcal F}$. Pour que $k\mathrel{-}k_\Phi$ soit surjectif, il \unsure{f} et il suffit que pour \deleted{toute constante, \unsure{les}} famille de parties \inserted{\uline{mutuellement disjointes}} $(E_\xi)_{\xi\in k}$ de $E$ indexée par $k$, \deleted{telle qu'on ait} \inserted{de réunion} $E$, \deleted{\unsure{$=\bigcup E_\xi$}} il existe un $\xi\in k$ tel que $E_\xi\in\Phi$. A fortiori, il \uline{faut} que $\Phi$ soit un ultrafiltre (puisque $\unsure{\mathrm{Card}|\mathrm{card}}\,k\geq2$ par hypothèse).

% archive sheet number: 599
